\documentclass[a4paper,11pt]{ctexart}

% 在子目录中编译时，通过相对路径寻找 mathnote-preamble
\InputIfFileExists{mathnote-preamble.tex}{}{%
  \InputIfFileExists{../mathnote-preamble.tex}{}{%
    \InputIfFileExists{../../mathnote-preamble.tex}{}{%
      \PackageError{mathnote-notes}{mathnote-preamble.tex 未找到}{请确认 mathnote-preamble.tex 与当前文稿的相对路径配置正确。}%
    }%
  }%
}

% \MathNoteEnablePrint % 如主要用于纸质打印，可取消注释

% 一些专题中用到的绘图宏包（如切线放缩与导数可视化）
\usepackage{pgfplots}
\pgfplotsset{compat=1.18}

% 告诉各单册笔记：当前处于“合订本模式”，不要再各自开启 document/class/目录
\newcommand{\MathNoteBookMode}{}

% 合订本元信息
\renewcommand{\notetitle}{高中数学笔记·合订本}
\renewcommand{\noteauthor}{（自填姓名）}
\renewcommand{\notedate}{\today}
\renewcommand{\notesubtitle}{集合·函数·导数·三角·数列·几何·概率}
\renewcommand{\noteversion}{v1.0}

% 在合订本中，把单册里的 section 降一级使用
\newcommand{\InputMathNoteAsSubsections}[1]{%
  \begingroup
    \let\MathNoteOrigSection\section
    \let\MathNoteOrigSubsection\subsection
    \let\MathNoteOrigSubsubsection\subsubsection
    \renewcommand{\section}{\subsection}
    \renewcommand{\subsection}{\subsubsection}
    \input{#1}%
  \endgroup
}

\begin{document}

% 封面
\thispagestyle{empty}
\PageTag[secondary]{高中数学·合订本}
\setcounter{page}{1}

\noindent
\begin{minipage}[t][0.7\textheight][c]{\textwidth}
  \raggedleft
  \vspace*{\fill}
  {\sffamily\bfseries\Huge\color{accent}\notetitle\par}
  \vspace{0.6em}
  {\Large\color{inkgray}\notesubtitle\par}
  \vspace{0.5em}
  {\large\color{inkgray}\noteauthor\par}
  \vspace{0.4em}
  \MathNoteVersionBlock
  \vspace*{\fill}
\end{minipage}

\clearpage
\thispagestyle{plain}
\phantomsection
\setcounter{tocdepth}{2}
\tableofcontents

\clearpage

% 1. 集合与不等式
\section{集合与不等式}
\InputMathNoteAsSubsections{集合与不等式.tex}

% 2. 函数与方程
\section{函数与方程}
\InputMathNoteAsSubsections{函数与方程.tex}

% 3. 导数与应用
\section{导数与应用}
\InputMathNoteAsSubsections{高中数学-导数与微积分.tex}

% 4. 三角函数与解三角形
\section{三角函数与解三角形}
\InputMathNoteAsSubsections{三角函数与解三角形.tex}

% 5. 平面向量与复数（占位，内容待补充）
\section{平面向量与复数}
本章内容尚在整理中，可在后续将对应专题笔记按与前几章相同的方式接入。

% 6. 数列
\section{数列}
\InputMathNoteAsSubsections{高中数学-数列复习笔记.tex}

% 7. 立体几何与空间向量（占位）
\section{立体几何与空间向量}
本章目前作为占位节，待你整理立体几何与空间向量的 tex 笔记后，可在此处通过
\verb|\InputMathNoteAsSubsections{...}| 接入。

% 8. 解析几何（占位）
\section{解析几何}
本章目前为空白结构，用于未来整合圆锥曲线、直线与圆、空间直角坐标系等解析几何专题。

% 9. 概率（占位）
\section{概率}
本章预留给离散与几何概率、古典概率模型以及随机变量的基础内容。

\end{document}
